% VI48-DETAILED-CHALLENGE-SOLUTIONS
% VI48-DETAILED-CHALLENGE-01
\subsection*{Challenge VI.48.1}
\begin{solution}
For a short exact sequence of cochain complexes, define the boundary by choosing a lift of a cocycle, applying the differential, and identifying the result with a cocycle in the kernel complex.  Verify independence of the lift, exactness at every term, and naturality by direct diagram chasing.  For sheaves, injective or flabby resolutions enter because applying global sections to the original sheaf sequence need not be right-exact; resolutions replace it by a short exact sequence of complexes on which the complex-level theorem applies.
\end{solution}

% VI48-DETAILED-CHALLENGE-02
\subsection*{Challenge VI.48.2}
\begin{solution}
Arrange the \(3\times3\) diagram as three compatible short exact rows and columns.  Applying cohomology produces commutative ladders by naturality of every connecting map.  Comparing the two ways of moving around each square yields the cohomological \(3\times3\) compatibility; at the resolution level it is the ordinary functoriality of boundary maps for a diagram of short exact complexes.
\end{solution}

% VI48-DETAILED-CHALLENGE-03
\subsection*{Challenge VI.48.3}
\begin{solution}
For
\[
0\to\mathcal L(-p)\to\mathcal L\to\mathcal L|_p\to0,
\]
the relevant part is
\[
0\to H^0(\mathcal L(-p))\to H^0(\mathcal L)\xrightarrow{\mathrm{ev}_p}k
\to H^1(\mathcal L(-p))\to H^1(\mathcal L)\to0.
\]
If evaluation is onto, \(h^0\) increases by one and \(h^1\) is unchanged when the point is added.  If evaluation is zero/not onto, the missing dimension reappears through the connecting map and changes \(h^1\).  Iterating this dichotomy is the elementary mechanism behind divisor-degree arguments and Riemann--Roch.
\end{solution}

% VI48-DETAILED-CHALLENGE-04
\subsection*{Challenge VI.48.4}
\begin{solution}
Choose a short exact sequence with a quotient section whose obstruction is nonzero on a coarse open.  The sheaf epimorphism still guarantees a lift near every point.  Cover the coarse open by such lifting neighborhoods and refine all finite intersections accordingly.  On the refined cover, the quotient map is surjective on every cochain term, giving a short exact sequence of Čech complexes.  This demonstrates concretely why fixed-cover termwise surjectivity is an extra hypothesis rather than part of the definition of a sheaf epimorphism.
\end{solution}

% VI48-DETAILED-CHALLENGE-05
\subsection*{Challenge VI.48.5}
\begin{solution}
Let \(0\to\mathcal F\to\mathcal I^0\to\mathcal Q^1\to0\) be the first step of a flabby resolution.  For \(n\ge2\),
\[
H^n(\mathcal F)\cong H^{n-1}(\mathcal Q^1).
\]
Repeat with \(0\to\mathcal Q^1\to\mathcal I^1\to\mathcal Q^2\to0\) until
\[
H^n(\mathcal F)\cong H^1(\mathcal Q^{n-1}).
\]
The last group is a global lifting cokernel, so a suitable surjectivity/vanishing argument in degree one proves higher vanishing.  This is dimension shifting as a practical proof strategy.
\end{solution}

